\section{Contoh Terhitung}
\label{sec:contoh-dh}

% Sumber: hlm. 15 dan 22.

Contoh pada bagian ini memakai bilangan yang sama dengan bahan kuliah. Setiap
angka dihitung ulang dengan program agar hasilnya dapat diperiksa.

\begin{example}[Pertukaran kunci dengan $p = 353$ dan $g = 3$]
\label{ex:dh-353}
Alice dan Bob menyepakati $p = 353$ dan $g = 3$. Kedua nilai itu memenuhi
syarat, yaitu $3 < 353$ dan $3$ adalah akar primitif dari $353$.

\begin{enumerate}
    \item Alice memilih $a = 97$ dan menghitung
          \[ A = g^a \bmod p = 3^{97} \bmod 353 = 40 . \]
          Alice mengirim $A$ kepada Bob.
    \item Bob memilih $b = 233$ dan menghitung
          \[ B = g^b \bmod p = 3^{233} \bmod 353 = 248 . \]
          Bob mengirim $B$ kepada Alice.
    \item Alice menghitung kunci rahasia $K$,
          \[ K = B^a \bmod p = 248^{97} \bmod 353 = 160 . \]
    \item Bob menghitung kunci rahasia $K$,
          \[ K = A^b \bmod p = 40^{233} \bmod 353 = 160 . \]
\end{enumerate}

Jadi, Alice dan Bob memperoleh kunci rahasia yang sama, yaitu $K = 160$.
Kunci itu dapat dipakai untuk enkripsi dan dekripsi pesan dengan algoritma
kunci-simetri.
\end{example}

\begin{example}[Parameter dari perkakas simulasi daring]
\label{ex:dh-1601}
Perkakas simulasi Diffie--Hellman meminta pengguna memasukkan dua bilangan,
$G$ dan $N$, dengan $N$ harus bilangan prima. Misalkan dipilih $G = 1601$ dan
$N = 4789$.

Bob memilih kunci privat $x = 20$ dan menghitung
\[ A = G^{x} \bmod N = 1601^{20} \bmod 4789 = 2716 . \]
Alice memilih kunci privat $y = 4$ dan menghitung
\[ B = G^{y} \bmod N = 1601^{4} \bmod 4789 = 1302 . \]

% K tidak dicantumkan pada hlm. 22. Nilainya dihitung dari prosedur yang sama
% dan sudah diperiksa dengan Python.
Setelah bertukar, keduanya menghitung kunci bersama:
\[
\begin{aligned}
K &= B^{x} \bmod N = 1302^{20} \bmod 4789 = 4257, \\
K &= A^{y} \bmod N = 2716^{4} \bmod 4789 = 4257 .
\end{aligned}
\]
Kunci bersama yang diperoleh adalah $K = 4257$.
\end{example}
